\chapter{Conclusions}
\label{chap:conclusions}

Local robotic polishing requires a strategy for movement, contact loading, and time within a selected region. Human region specification supplies the spatial task, but does not by itself determine that temporal strategy. This thesis has addressed the resulting skill-generation problem through paired motion--force teaching, observation-conditioned sequence learning, and a common robot execution interface.

Tool position is learned relative to a frozen stain center and paired with orientation and target force on a shared sequence. The direction-transfer adapter aligns live XY and the policy image with the teaching direction, then restores generated XY by inverse rotation and center addition. Z, tool orientation, and force keep their conventions. The evaluated Flow Matching policy uses RGB, Cartesian state, and three-axis force history, with explicit coordinate restoration, pose conditioning, plan connection, and existing contact regulation. E1/E2 disable direction alignment and do not measure its transfer performance. The implementation separates upper-level skill generation, lower-level contact control, and completion recognition.

The first research question compares strategy providers. In the recorded E1 nominal cohort, rules, teaching replay, and the conditioned policy all show the required loaded ordered pass. The original completion counts differ, but the height-reference sensitivity analysis shows that the difference cannot be interpreted directly as a physical exit advantage for the policy. Replay provides a fast, repeatable nominal strategy, and the conditioned policy has greater excess-force exposure in this comparison. The implemented system therefore supports learning as an executable option alongside strong fixed alternatives; the current results do not establish its overall superiority. Performance under actual specimen-height changes remains unevaluated.

The second research question separates force output learning from force observation. E2 exercises all three implemented contracts while retaining lower-level feedback. The force-observing condition has shorter recorded operator-based completion intervals in this cohort, while A4 in the motion-only condition fails to progress to release and is interrupted. That attempt remains in the condition's five-attempt force and timing statistics, with successful completion time defined only for the other four attempts. These results reveal differences in phase progression and execution dispersion. The ordered, small, single-seed comparison does not isolate a universal causal advantage of either force-information choice.

The demonstrated scope is the generation and execution of a human-guided local contact skill on the marked-surface setup. It does not include autonomous recognition of defect severity, learned recognition of complete scratch removal, calibrated material-removal estimation, or measured surface restoration. The existing dose, blending, and adaptive-execution assets provide concrete extension paths, with their implementation and validation status stated separately from the main findings.

The thesis's contribution is the coherent connection from local task specification to a time-related motion--force representation, conditioned generation, shared execution, and corresponding evaluation. That connection provides a basis for deciding when rules or replay suffice and for testing, with additional physical conditions and quality measurements, when learned generation adds useful capability.
